% chapter/chapter4.tex
% Beasts of the Wilds – Teeth, Claws, and the Hunger That Does Not Bargain
% Rewritten in the voice of Lerris Fenwood, with marginalia from the Gray Wanderer,
% Saikou Ira, and Dusana of the Raven Road.

\epigraph{“The wilds do not hate you. They do not love you. They simply remember what you owe.”}{---Dusana of the Raven Road, from a fire‑side tale in the Violet Steppe}

The wilds are not a monster. They are a ledger written in teeth, frost, and the silence after a scream. I learned this on the road to Ash‑Fenn, when a wolf followed our caravan for three days. It did not attack. It waited. On the fourth day, the caravan master paid a debt he had forgotten. The wolf turned and walked away. The wilds do not forget. They do not forgive. But they will sometimes accept payment.

This chapter catalogs the beasts that walk the spaces between towns: predators of the steppe, the forest, the coast, and the deep places beneath the world. Some are natural; most are not. All are hungry.

\section*{Steppe & Grasslands}

\subsection*{Earth‑Shark}

\noindent\textbf{Description}: Burrows beneath loess and tundra; dorsal ridge like a plowshare. They are not sharks. They are something older that learned the shape of hunger from the sea and brought it to land.

\noindent\textbf{Stats}: Tier III, Cap 4. Corporeal, Burrowing, Ambush.\\
\textbf{Harm}: 5. \textbf{Fatigue}: Body 4.

\noindent\textbf{Traits}: 
\emph{Breach} (DV 3, Harm 2) – erupts from the ground beneath a target, knocking them prone. \emph{Tremorsense} – detects any creature touching the ground within Near range. Cannot be surprised from below.

\noindent\textbf{Complications}: The Earth‑Shark will not attack a group larger than three. It prefers isolated prey. The furrows it leaves breathe; they pulse with a slow, warm exhalation.

\noindent\textbf{Resolution}: Lure with drumbeats or goat bells. The vibration confuses its tremorsense. A silver blade driven into the ridge between its eyes will kill it instantly; otherwise, fire drives it back underground.

\noindent\textbf{Clock}: \emph{Tremor [6]} – ticks each time the Earth‑Shark breaches. When full, the ground becomes unstable; all actions at Desperate Position.

\gray{I have seen an Earth‑Shark breach beneath a horse. The horse did not scream. The shark pulled it under. The ground closed. There was no blood.}

\subsection*{Dust Hulk}

\noindent\textbf{Description}: Wind‑packed silhouette that rises from dunes. It is not a creature; it is a \textbf{memory} of a creature, given weight by the desert’s thirst.

\noindent\textbf{Stats}: Tier II, Cap 3. Corporeal (semi‑incorporeal), Wind‑Bound, Silent.\\
\textbf{Harm}: 4. \textbf{Fatigue}: Spirit.

\noindent\textbf{Traits}: 
\emph{Dust Cloud} (DV 2, Fatigue 1) – all within Near range suffer −1 die to sight‑based actions. \emph{Silent Move} – makes no sound; cannot be detected by hearing.

\noindent\textbf{Complications}: The Dust Hulk cannot cross running water. It will follow travelers for days, never attacking, only watching. Its presence is announced by a sudden stillness in the wind.

\noindent\textbf{Resolution}: Spill water on its shadow. The shadow will dissolve, and the Hulk will collapse into a pile of harmless dust. A single drop is enough; it must be fresh, drawn from a well or spring.

\noindent\textbf{Clock}: \emph{Thirst [6]} – ticks each time the Hulk is denied water. When full, it becomes a \textbf{Dust Walker} (Tier IV) – corporeal, aggressive, and able to drain the moisture from living flesh.

\saikou{I encountered a Dust Hulk on the road to Midh Ahkaz. I poured a canteen of water onto its shadow. The shadow drank. The Hulk sighed and became sand. The sand was warm. I kept a pinch. It is still warm.}

\subsection*{Bone Kite}

\noindent\textbf{Description}: Carrion bird with rune‑etched pinions. Its feathers are not feathers; they are thin slivers of bone. It does not fly; it drifts.

\noindent\textbf{Stats}: Tier I, Cap 2. Corporeal, Aerial, Opportunistic.\\
\textbf{Harm}: 2. \textbf{Fatigue}: Body 1.

\noindent\textbf{Traits}: 
\emph{Scavenger’s Bite} (DV 2, Harm 1) – prefers to attack the wounded. \emph{Rune Sight} – can see the weight of a creature’s unfulfilled promises. A creature with an active Debt Timer is visible from twice the normal range.

\noindent\textbf{Complications}: Bone Kites do not attack healthy prey. They circle where a death is expected. Their runes change with the weather; reading them requires a Wits + Lore test (DV 4).

\noindent\textbf{Resolution}: Offer a name‑stake – a twig wrapped with a lock of hair – as a barter. The Kite will take the stake and leave. It will remember the name and may return with a message from the dead.

\noindent\textbf{Clock}: \emph{Scavenger’s Patience [4]} – ticks each time the Kite is refused an offering. When full, it summons a flock (2d6 Bone Kites).

\dusana{The Bone Kites spoke to me once. They said a name I had not heard in twenty years. I paid them with a story. The story was about a woman who forgot to bury her husband. The kites listened. They did not judge. They simply remembered.}

\subsection*{Salt Screamer}

\noindent\textbf{Description}: Pale eel in dry salt pans; voices like flutes. They are not eels. They are the \textbf{nerves} of the desert, exposed and singing.

\noindent\textbf{Stats}: Tier II, Cap 2. Corporeal, Subterranean, Musical.\\
\textbf{Harm}: 3. \textbf{Fatigue}: Body 2.

\noindent\textbf{Traits}: 
\emph{Scream} (DV 3, Fatigue 1) – all within Near range must test Resolve (DV 3) or be Stunned for one exchange. \emph{Salt Glide} – moves through salt pans as if swimming; cannot be tracked.

\noindent\textbf{Complications}: The Salt Screamer’s scream is not an attack; it is a \textbf{question}. The question is always: “What is your name?” Answering truthfully grants +1 die to all actions for the scene. Answering falsely advances the \emph{Scream’s Patience} timer.

\noindent\textbf{Resolution}: Scatter ash. The Screamer will coil and sing to forget. It will not attack for one scene. To kill it, you must pour fresh water into its burrow; the Screamer will dissolve, leaving only a salt crystal the size of a thumb.

\noindent\textbf{Clock}: \emph{Scream’s Patience [4]} – ticks each time a lie is answered. When full, the Screamer’s scream becomes Harm 2 and affects all within Far range.

\lerris{I kept a Salt Screamer’s crystal on my desk for a year. It sang at midnight. The song was always a name I had tried to forget. I threw the crystal into the river. The river sang back. I do not know the tune.}

\section*{Forests & Jungles}

\subsection*{Phase Panther}

\noindent\textbf{Description}: Shadow‑striped cat that steps between gusts. Its stripes are not fur; they are \textbf{cracks} in the air, where the panther has learned to fold itself.

\noindent\textbf{Stats}: Tier III, Cap 4. Corporeal, Phasing, Silent.\\
\textbf{Harm}: 5. \textbf{Fatigue}: Body 3.

\noindent\textbf{Traits}: 
\emph{Phase Step} – once per scene, the panther may teleport to any location within Near range, ignoring opportunity attacks. \emph{Shadow Stripes} – while in dim light or darkness, the panther has Dominant Position on Stealth rolls.

\noindent\textbf{Complications}: The panther cannot phase through mirrored water. A pool of still water, even a small one, acts as a barrier. It also cannot phase while bleeding; a wound that draws blood will hold it in place for one exchange.

\noindent\textbf{Resolution}: Mirrored water. A silver mirror held toward the panther will reflect its own stripes back at it; the panther will become fascinated (Stunned for one exchange). A silver blade through the heart kills it instantly.

\noindent\textbf{Clock}: \emph{Hunger [6]} – ticks each time the panther phases. When full, it phases into a creature’s shadow and emerges inside their body. The victim takes Harm 3 and is Incapacitated.

\gray{I have seen a Phase Panther’s shadow without its body. The shadow was hunting alone. It had forgotten its flesh.}

\subsection*{Lichen‑Bear}

\noindent\textbf{Description}: Moss‑matted ursine with rock‑teeth. The moss is not a covering; it is the bear’s \textbf{memory} of a forest that no longer exists.

\noindent\textbf{Stats}: Tier III, Cap 5. Corporeal, Sturdy, Ancient.\\
\textbf{Harm}: 7. \textbf{Fatigue}: Body 5.

\noindent\textbf{Traits}: 
\emph{Maul} (DV 3, Harm 3). \emph{Stone Heart} – ignores the first Harm 2 each scene. \emph{Moss Breath} – once per scene, exhale a cloud of spores; all within Close range test Endurance (DV 4) or gain the Condition \emph{Spored} (the bear knows your location for the rest of the scene).

\noindent\textbf{Complications}: The Lichen‑Bear is slow but relentless. It does not run. It walks. It will follow for days.

\noindent\textbf{Resolution}: Smoke of juniper calms it; salt enrages it. To kill it, you must burn the moss. A torch held to its back will deal +2 Harm, but the bear will frenzy for one exchange. The rock‑teeth can be harvested; they are cold and never warm.

\noindent\textbf{Clock}: \emph{Forest’s Memory [8]} – ticks each time the bear is wounded. When full, the moss grows over the bear’s eyes, and it becomes blind but berserk (+2 dice to attack, −2 to defense DV).

\saikou{I tracked a Lichen‑Bear for a week. It led me to a grove where a woman had been buried alive. The moss had grown over her grave. The bear was not a predator. It was a warden.}

\subsection*{Bramble Knight}

\noindent\textbf{Description}: Empty helm grown through with thorn. It is not a knight. It is a \textbf{hedge} that learned to walk, still dreaming of the wall it once was.

\noindent\textbf{Stats}: Tier II, Cap 3. Corporeal, Animated, Territorial.\\
\textbf{Harm}: 4. \textbf{Fatigue}: Body 3.

\noindent\textbf{Traits}: 
\emph{Thorn Lash} (DV 2, Harm 1 + 1 Fatigue from bleeding). \emph{Iron Fear} – if presented with naked iron, the Knight must test Resolve (DV 3) or lose its next action. \emph{Plant Mind} – immune to fear, charm, and mind‑affecting effects.

\noindent\textbf{Complications}: The Bramble Knight cannot cross a line of salt. It will patrol the same boundary for years, never straying. Its helm is always empty; the knight is the thorns, not the metal.

\noindent\textbf{Resolution}: Present a hedge‑blessing – a branch of hawthorn tied with red thread. The Knight will escort you through its territory without attacking. To destroy it, burn the thorns. The helm will become a simple, rusted helm. It may be worn; the wearer gains +1 die to resist fear, but the thorns will grow slowly into their skin.

\noindent\textbf{Clock}: \emph{Hedge’s Patience [6]} – ticks each time the Knight is denied a boundary. When full, the Knight summons 1d4 lesser Bramble Knights (Cap 1, Harm 2).

\dusana{The Bramble Knight in the Valewood wore a helm that had belonged to a king. The king had been dead for three hundred years. The thorns did not care. They only remembered the border.}

\section*{Coasts & Swamps}

\subsection*{Reef Serpent}

\noindent\textbf{Description}: Kelp‑maned constrictor nesting in pilings. It is not a snake; it is a \textbf{rope} that forgot it was tied.

\noindent\textbf{Stats}: Tier II, Cap 3. Corporeal, Amphibious, Constrictor.\\
\textbf{Harm}: 4. \textbf{Fatigue}: Body 3.

\noindent\textbf{Traits}: 
\emph{Constrict} (DV 3, Harm 1 + target cannot move until freed; escape requires Body + Athletics vs DV 4). \emph{Kelp Mane} – provides partial cover (+1 Armor) against ranged attacks.

\noindent\textbf{Complications}: The Reef Serpent is territorial but not aggressive. It will warn intruders by drumming its tail against the pilings. The drumming is in eights. If you hear a ninth drum, it is already behind you.

\noindent\textbf{Resolution}: Beat a drum in a rhythm the serpent does not recognize. It will become confused and withdraw. To kill it, you must sever the kelp mane; without it, the serpent loses its cover and becomes Vulnerable (−1 to defense DV).

\noindent\textbf{Clock}: \emph{Tide’s Hunger [6]} – ticks each time the serpent constricts a victim. When full, it drags the victim underwater and drowns them (Harm 3).

\lerris{I saw a Reef Serpent coiled around a piling. It had been there for a century. The kelp had grown into the wood. The serpent was not waiting for prey. It was waiting for the piling to remember it was a tree.}

\subsection*{Grave‑Eel}

\noindent\textbf{Description}: Burrows through barrow loam; tastes names. It is not an eel. It is a \textbf{question} that learned to wiggle.

\noindent\textbf{Stats}: Tier I, Cap 2. Corporeal, Burrowing, Curious.\\
\textbf{Harm}: 2. \textbf{Fatigue}: Body 1.

\noindent\textbf{Traits}: 
\emph{Name Taste} – if the Grave‑Eel touches a creature, the GM learns one true name that creature has been called. \emph{Sinking burrow} – the eel can collapse the ground beneath it; creatures within Close must test Body + Athletics (DV 3) or be knocked prone.

\noindent\textbf{Complications}: Grave‑Eels are not aggressive. They are \textbf{scholars}. They will follow a party for hours, observing, tasting names from the air.

\noindent\textbf{Resolution}: Feed it a false name. The eel will become confused and retreat. To kill it, you must sing a lineage – a genealogy of at least three generations. The eel will listen, then burrow away. It is not harmed; it is satisfied.

\noindent\textbf{Clock}: \emph{Name Hunger [6]} – ticks each time the eel tastes a true name. When full, it surfaces and speaks a name aloud. The name is always a secret someone wanted to keep.

\saikou{I fed a Grave‑Eel a false name. The name was “Ira.” The eel tasted it and spat. It knew. They always know.}

\section*{Mountains & Underways}

\subsection*{Vein‑Serpent (Aeler)}

\noindent\textbf{Description}: Serpent of polished copper and malachite, slithering through stone as if it were water. It is not a serpent. It is the mountain’s \textbf{blood} given form.

\noindent\textbf{Stats}: Tier II, Cap 3. Corporeal, Stone‑Swim, Mineral.\\
\textbf{Harm}: 4. \textbf{Fatigue}: Body 3.

\noindent\textbf{Traits}: 
\emph{Stone Swim} – moves through solid rock as if it were water; cannot be targeted while submerged. \emph{Copper Bite} (DV 2, Harm 1 + 1 Fatigue; the bite is cold, not hot). \emph{Mineral Mind} – immune to poison, disease, and mind‑affecting effects.

\noindent\textbf{Complications}: The Vein‑Serpent is not hostile unless the stone is disturbed. Mining, drilling, or breaking a keystone will draw its attention.

\noindent\textbf{Resolution}: Negotiate. The serpent will accept a promise to maintain the stability of the stone – for every ounce of ore taken, an ounce of mortar must be laid. To kill it, you must collapse the tunnel. The serpent will die, but so will the miners.

\noindent\textbf{Clock}: \emph{Stone’s Grudge [6]} – ticks each time a keystone is broken without repair. When full, the serpent collapses a tunnel on the offenders.

\gray{The Aeler do not kill Vein‑Serpents. They hire them. The serpent is paid in copper dust and silence.}

\subsection*{Wind Drake}

\noindent\textbf{Description}: Rib‑thin sky serpent that rides fronts. Its bones are hollow; its scales are feathers made of glass.

\noindent\textbf{Stats}: Tier III, Cap 4. Corporeal, Aerial, Storm‑Born.\\
\textbf{Harm}: 5. \textbf{Fatigue}: Spirit.

\noindent\textbf{Traits}: 
\emph{Lightning Breath} (DV 4, Harm 2 + 1 Fatigue from shock). \emph{Storm Rider} – while in a storm, the drake has Dominant Position on all rolls. \emph{Feathers of Glass} – ranged attacks against the drake suffer −1 die (the feathers deflect).

\noindent\textbf{Complications}: The Wind Drake is territorial during storms. It will attack anyone who flies a banner or flag – it mistakes the cloth for a rival drake.

\noindent\textbf{Resolution}: Whistle back the storm. A successful Wits + Survival test (DV 5) can calm the weather; the drake will lose interest and fly away. To kill it, you must ground it – a net weighted with copper will hold it for one exchange. Silver through the heart kills it instantly.

\noindent\textbf{Clock}: \emph{Storm’s Fury [8]} – ticks each time the drake breathes lightning. When full, it summons a second drake.

\dusana{I saw a Wind Drake fall from the sky. It had been struck by a bell’s note. The note was nine. The drake had not learned to stop counting.}

\section*{Dungeon Living Infrastructure}

\subsection*{Rust Moth}

\noindent\textbf{Description}: Swarm of iron‑scaled moths that taste memory. They are not insects; they are the dungeon’s \textbf{forgetting} given wings.

\noindent\textbf{Stats}: Tier I, Cap 2 (swarm). Corporeal, Swarm, Rusting.\\
\textbf{Harm}: 2 (area). \textbf{Fatigue}: Body.

\noindent\textbf{Traits}: 
\emph{Rust Touch} – any metal object touched by the swarm becomes Compromised after one exchange, Neglected after two, and destroyed after three. \emph{Memory Feed} – the swarm consumes memories as rust consumes iron. A character who takes Harm from the swarm must test Spirit + Resolve (DV 3) or forget a minor fact.

\noindent\textbf{Complications}: The swarm is attracted to old keys. A character carrying a key that has not been used in a decade will be targeted first.

\noindent\textbf{Resolution}: Feed them old keys. The swarm will ignore living creatures for one scene. To destroy the swarm, ring a copper bell; the sound repels them. A single bell‑ring will drive them back to the walls.

\noindent\textbf{Clock}: \emph{Rust Memory [6]} – ticks each time the swarm consumes a memory. When full, it becomes a \textbf{Rust Walker} (Tier III, Cap 4) – a humanoid figure of rust that can wield corroded weapons.

\lerris{I fed a Rust Moth swarm a key that had opened my father’s study. The key was warm when I gave it. The swarm took it and left. The key is still warm. I do not go into the study.}

\subsection*{Lantern‑Eye}

\noindent\textbf{Description}: Floating orb of watchlight; jealous of secrets. It is not a creature; it is the dungeon’s \textbf{attention} made visible.

\noindent\textbf{Stats}: Tier II, Cap 3. Incorporeal, Floating, Inquisitive.\\
\textbf{Harm}: Cannot be harmed by mundane weapons. \textbf{Fatigue}: Spirit.

\noindent\textbf{Traits}: 
\emph{Watchlight} – the Lantern‑Eye illuminates a zone; stealth within that zone is impossible. \emph{Secret Jealousy} – if a character speaks a secret aloud within sight of the Eye, the Eye will become hostile and attempt to blind them (DV 3, Fatigue 1).

\noindent\textbf{Complications}: The Lantern‑Eye is not malevolent. It is \textbf{curious}. It will follow the party, watching, until a secret is spoken.

\noindent\textbf{Resolution}: Offer a whispered truth. The Eye will accept the truth and part, revealing a hidden passage. To destroy it, you must throw a shard of obsidian into its light. The light will go out. The Eye will fall and shatter.

\noindent\textbf{Clock}: \emph{Watchlight’s Patience [6]} – ticks each time a secret is kept from the Eye. When full, the Eye summons 2d4 lesser Lantern‑Eyes (Cap 1, Harm 1 each).

\saikou{I offered a Lantern‑Eye a truth: a confession of a murder I did not commit. The Eye accepted. It showed me the passage to the murderer’s cell. The murderer was already dead. The Eye did not care. It only wanted a secret.}

\section*{Quick Reference – Beasts of the Wilds}

\begin{tabular}{|l|l|l|l|}
\hline
\textbf{Creature} & \textbf{Tier} & \textbf{Vulnerability} & \textbf{Resolution} \\
\hline
Earth‑Shark & III & Silver, fire & Kill or drive off with drums \\
Dust Hulk & II & Water & Pour water on shadow \\
Bone Kite & I & Name‑stake & Barter with name‑stake \\
Salt Screamer & II & Fresh water, ash & Kill with water or quiet with ash \\
Phase Panther & III & Mirrored water, silver & Kill or fascinate with mirror \\
Lichen‑Bear & III & Fire, juniper smoke & Kill with fire or calm with smoke \\
Bramble Knight & II & Salt, fire & Hedge‑blessing or burn \\
Reef Serpent & II & Sever kelp mane & Kill or confuse with drum \\
Grave‑Eel & I & False name & Feed false name or sing genealogy \\
Vein‑Serpent & II & Keystone clause & Negotiate or collapse tunnel \\
Wind Drake & III & Copper net, silver & Ground or kill with silver \\
Rust Moth & I & Copper bell, old keys & Feed keys or repel with bell \\
Lantern‑Eye & II & Obsidian shard, truth & Offer truth or destroy with obsidian \\
\hline
\end{tabular}

\lerris{I keep this table on a loose leaf. The leaf is pinned to the inside cover of my ledger. The pin is copper. The ledger does not rust.}

\gray{The wilds do not hate you. They do not love you. They simply remember. And memory is a debt that never closes.}
\saikou{I have killed seven Earth‑Sharks. Each time, the ground was warm for a week. The grass grew back thicker. The sharks did not.}
\dusana{The Bone Kites followed me for a year after I told them the story of the woman who forgot to bury her husband. They remembered. They always remember.}

